% Generated draft: requires mathematical and source review.
\documentclass[11pt]{article}
\usepackage[margin=25mm]{geometry}
\usepackage{fontspec}
\setmainfont{Latin Modern Roman}
\usepackage{amsmath,amssymb,mathtools,unicode-math}
\setmathfont{Latin Modern Math}
\usepackage{enumitem,needspace,xurl,hyperref}
\setlength{\parindent}{0pt}
\setlength{\parskip}{5pt}
\setlength{\emergencystretch}{3em}
\newcommand{\N}{\mathbb{N}}
\newcommand{\Z}{\mathbb{Z}}
\newcommand{\Q}{\mathbb{Q}}
\newcommand{\R}{\mathbb{R}}
\newcommand{\C}{\mathbb{C}}
\newcommand{\F}{\mathbb{F}}
\newcommand{\A}{\mathbb{A}}
\newcommand{\PP}{\mathbb{P}}
\newcommand{\E}{\mathbb{E}}
\newcommand{\eps}{\varepsilon}
\newcommand{\dd}{\,\mathrm{d}}
\begin{document}
\section*{UnsolvedMath ID: 30006019}
\textbf{Source problem code:} catalogue entry
\subsection{Definitions and mathematical statement}
For a compact Riemann surface \(X\) of genus \(g\), a meromorphic \(k\)-differential is a nonzero meromorphic section of \(K_X^{\otimes k}\), where \(K_X\) is the canonical bundle and \(k\geq 1\) is an integer. Abelian and quadratic differentials correspond to \(k=1\) and \(k=2\), respectively. Write \(\Omega^{(k)}(\mu)\) for the stratum denoted in the source, with prescribed orders \(\mu\) of zeros and poles. If \(\mu=(m_1,\ldots,m_n)\) lists all such orders, then \(\sum_i m_i=k(2g-2)\). The supplied record does not specify the allowed pole orders, marking conventions, connected components, or whether differentials that are powers of lower-order differentials are included; these conventions remain unresolved.

For a fixed stratum, let \(\mathcal T_{k,\mu}(N)\) denote the finite collection described in the source as square-tiled surfaces in that stratum with at most \(N\) squares. Whenever this collection is nonempty, equip it with its uniform probability measure. The precise equivalence and marking conventions for this collection are not supplied. For an Abelian square-tiled surface \(S\), the horizontal and vertical directions give periodic cylinder decompositions. If \(\mathcal C_d(S)\) is the collection of cylinders in direction \(d\in\{\mathrm{horiz},\mathrm{vert}\}\), define the normalized cylinder-area multiset by
\[
 A_d(S)=\left\{\!\left\{
 \frac{\operatorname{Area}(C)}{\operatorname{Area}(S)}:
 C\in\mathcal C_d(S)
 \right\}\!\right\}.
\]
Multiplicity is retained, every entry is positive, and the entries sum to one. The source gives an example with \(A_{\mathrm{horiz}}(S)=A_{\mathrm{vert}}(S)=\{\!\{1/3,2/3\}\!\}\). These periodicity assertions concern square-tiled surfaces; periodicity is not asserted here for arbitrary Abelian differentials.

Let
\[
 \mathscr A=\left\{\{\!\{a_1,\ldots,a_r\}\!\}:
 r\geq1,\quad a_i>0,\quad \sum_{i=1}^r a_i=1\right\}.
\]
The source states that cylinder constructions for quadratic and quartic differentials can be obtained using a covering construction. It does not specify that construction or its area-counting conventions. Subject to those source conventions, the supplied fixed-stratum theorem asserts, for Abelian and quadratic strata,
\[
 \operatorname{Law}_{S\in\mathcal T_{k,\mu}(N)}
 \bigl(A_{\mathrm{horiz}}(S),A_{\mathrm{vert}}(S)\bigr)
 \longrightarrow \nu_{k,\mu}\otimes\nu_{k,\mu}
 \qquad (N\to\infty),\qquad k\in\{1,2\}.
\]
Thus the two limiting marginals agree and are independent. The topology or convergence convention on multisets is not specified in the supplied excerpt. Its footnote states that weighting each surface by \(1/|\operatorname{Aut}(S)|\) is more natural and gives the same limit; this is a supplied assertion, not a verification in this draft. The excerpt describes an analogous result for quartic differentials as work in progress and says that only a single cylinder-area vector \(A(S)\) is considered there. It does not supply the precise quartic theorem.

\textbf{Requested problem.} Determine the large-genus asymptotic behavior of the fixed-stratum limiting cylinder-area measures. More explicitly, for each fixed differential order \(k\) covered by the intended formulation and each admissible family of strata \(\Omega^{(k)}(\mu_g)\) of genus \(g\), determine the asymptotic behavior of
\[
 \nu_{k,\mu_g}\qquad\text{as }g\to\infty,
\]
with the counting limit \(N\to\infty\) taken first, separately in every fixed stratum. This is an iterated-limit problem, not a specified simultaneous limit in \(N\) and \(g\). Where the requisite fixed-stratum measure has not been defined or established in the supplied evidence, its construction and existence are prerequisites to this formulation.

The source says that a complete answer was then available only for principal quadratic strata, displayed with order data \((1^{4g-4})\), meaning \(4g-4\) simple zeros in the usual holomorphic principal-stratum interpretation for \(g\geq2\). No explicit principal-stratum limiting law is supplied. The cleaned statement requests all Abelian, quadratic, and higher-order strata, whereas the original excerpt discusses Abelian and quadratic strata and work in progress for quartic differentials. Consequently the all-orders scope is a supplied extension whose cylinder observable and counting model are not defined for general \(k\). The admissible families \(\mu_g\), any uniformity over strata, the limiting state space, the mode of convergence, and any necessary rescaling or reordering are also unspecified. These omissions prevent extracting a unique fully quantified asymptotic conjecture; the statement above records the research target without choosing those missing data.
\subsection{Short English statement}
First take the number of tiles to infinity within a fixed stratum to obtain a probability law for normalized cylinder areas, and then study how that law behaves as the genus grows. The supplied record reports a fixed-stratum theorem for Abelian and quadratic differentials, a known principal-quadratic-stratum answer, and quartic work in progress. Its broader all-strata, all-orders target needs additional conventions and a precise notion of large-genus convergence.
\subsection{Sources}
\begin{itemize}
\item Local UnsolvedMath registry: \texttt{data/unsolvedmath/problems.json}, permanent ID \(30006019\), source problem code \texttt{catalogue entry}, title \emph{Large-Genus Cylinder Decompositions of Square-Tiled Surfaces}. This draft uses the supplied complete record; it does not claim inspection of the registry file. No individual proposer is named in the record.
\item Supplied report citation: \emph{Riemann Surfaces: Random, Flat, and Hyperbolic Geometry} (2025), Oberwolfach Reports, pp.\ 2369--2424, DOI \texttt{10.4171/owr/2024/41}: \url{https://doi.org/10.4171/owr/2024/41}. The excerpt identifies the discussion as ``Problem 1: cylinder decompositions in large genus'' and includes the page marker 2387. The record gives proposed year 2024. The report citation, DOI year, and proposed year are retained as supplied without resolving their bibliographic relationship.
\item Supplied literature source labeled Delecroix--Goujard--Zograf--Zorich: \url{https://arxiv.org/abs/2007.04740}. The supplied excerpt attributes its fixed-stratum theorem to reference [1], but does not provide that reference's bibliographic entry; identification of [1] with this literature link is not assumed.
\item Supplied literature assessment dated exactly \texttt{2026-08-22}: the fixed-stratum counting limit is reported as proved, principal quadratic large-genus behavior as understood, and the all-strata and quartic program as incomplete. The registry labels the problem \texttt{partially\_solved}. This is supplied dated curation, not an independently verified assessment of the literature. The record itself requests independent verification of time-sensitive status and cited claims.
\end{itemize}
\end{document}
